\begin{tikzpicture}[font=\small,>=stealth,line join=round]
\draw[gray!45,rounded corners=3pt] (0,0) rectangle (14.2,9.4);
\node at (7.1,8.8) {约去零因子后，两条不同曲线逼近同一个高度};
\begin{scope}[shift={(3.8,2.8)},x=4cm,y=6cm]
\draw[->,gray] (-.65,0)--(.7,0) node[above] {$t=x-1$};
\draw[->,gray] (0,-.1)--(0,.78) node[left] {$y$};
\draw[gray,dotted] (-.6,{1/3})--(.6,{1/3});
\draw[blue!75!black,thick,domain=-.5:.5,samples=80] plot (\x,{1/(3+2*\x)});
\draw[orange!85!black,thick,dashed,domain=-.5:.5,samples=80] plot (\x,{1/(3+3*\x)});
\draw[fill=white] (0,{1/3}) circle(2pt);
\node[above left=14pt] at (0,{1/3}) {$\frac13$};
\end{scope}
\node[anchor=west,align=left] at (8.8,5.1)
{蓝：原比值约分后\\[8pt]$\dfrac1{3+2t}$\\[17pt]
橙：求导比值约分后\\[8pt]$\dfrac1{3+3t}$\\[17pt]
$t\to0$ 时都趋于 $\frac13$。};
\node[text width=13.2cm,align=center] at (7.1,.65) {空心点提醒原式在 $x=1$ 处是 $0/0$；极限研究的是靠近该点时的数值。};
\end{tikzpicture}
